\chapter{Álgebra genealógica de caminos y representación de Feynman}
\label{chap:integral-genealogica}

La formulación de Feynman asigna a cada historia posible una amplitud y suma
esas amplitudes para obtener el propagador. En HMT la palabra \emph{historia}
posee un contenido anterior a la amplitud: una ruta conserva posición, hoja,
orientación, fase, acarreo, frontera, memoria y condición de supervivencia. La
suma de caminos no se añade desde fuera a ese sistema. Aparece al representar
linealmente su categoría de rutas
\cite{Feynman1948,FeynmanHibbs1965}.

El desarrollo posee tres niveles matemáticos enlazados. La expansión de una
potencia matricial como suma finita de caminos ya está demostrada en Mecánica
Dimensional. El ledger \(27\times27\) proporciona una realización unitaria
con shift, mezclador y fase, diagonalizable por Fourier. La familia de sumas
finitas admite después un límite cilíndrico cuando la medida, la acción y el
refinamiento satisfacen las hipótesis del teorema de convergencia. Así quedan
reunidas la identidad combinatoria, su realización espectral y su extensión
funcional.

\section{La categoría enriquecida de rutas}

Sea \(X\) un conjunto finito de estados observables y sea
\(\mathsf P_{\TPK}\) la categoría libre generada por las transiciones
admisibles del TPK. Denotemos
\(\widetilde X=\operatorname{Ob}(\mathsf P_{\TPK})\). Sus objetos son estados
enriquecidos de frontera. Un morfismo
\(\gamma:\widetilde x\to\widetilde y\) es una secuencia composable de
transiciones enriquecidas.
La composición es concatenación y la identidad es el camino vacío.

La proyección visible
\[
 q_{\rm vis}:\mathsf P_{\TPK}\longrightarrow \mathsf P_X
\]
olvida parte de la fibra. Dos morfismos diferentes pueden tener la misma
palabra observable. Esta no fidelidad es esencial: una suma que agregase antes
de tiempo sobre \(\mathsf P_X\) perdería precisamente la información capaz de
distinguir futuros incompatibles.

\begin{definition}[Peso genealógico]
Sea \(\mathcal B\) un álgebra real normada, no necesariamente conmutativa, y
sea \(B(\mathcal B,\cdot,1)\) la categoría de un objeto cuyo monoide de
endomorfismos es \((\mathcal B,\cdot,1)\). Un peso genealógico es un funtor
\[
 W:\mathsf P_{\TPK}\longrightarrow B(\mathcal B,\cdot,1)
\]
que satisface
\[
 W(\id_x)=1_{\mathcal B},\qquad
 W(\gamma_2\circ\gamma_1)=W(\gamma_2)W(\gamma_1).
\]
El orden del producto sigue el orden causal de la ruta.
\end{definition}

El peso puede contener un transporte de fibra, una penalización positiva, una
acción, un carácter del ledger o varios de esos factores. La definición no
confunde esas funciones: cada componente mantiene su dominio y su regla de
composición.

\section{La suma finita exacta}

Sea \(E_x\) la fibra lineal sobre \(x\in X\) y
\(\mathcal H_X=\bigoplus_{x\in X}E_x\). Sea
\(U=(U_{yx})_{x,y\in X}\) el operador global por bloques, con
\(U_{yx}:E_x\to E_y\). Interpretamos \(U_{yx}\) como el transporte elemental
desde \(x\) hasta \(y\). Para un camino visible
\(\gamma=(x_0,x_1,\ldots,x_R)\), definimos
\[
 W_U(\gamma)=U_{x_Rx_{R-1}}\cdots U_{x_1x_0}.
\]

\begin{theorem}[Expansión finita de caminos]
\label{thm:suma-finita-caminos}
Para todo \(R\geq 1\),
\begin{equation}
 (U^R)_{fi}
 =\sum_{\substack{\gamma:i\to f\\|\gamma|=R}}W_U(\gamma).
 \label{eq:suma-finita-caminos}
\end{equation}
Si el operador global \(U:\mathcal H_X\to\mathcal H_X\) es ortogonal ---o
unitario tras elegir una carta compleja---, la
evolución \(\psi_R=U^R\psi_0\) conserva la norma.
\end{theorem}

\begin{proof}
La entrada \((f,i)\) del producto \(U^R\) es
\[
 \sum_{x_1,\ldots,x_{R-1}}
 U_{fx_{R-1}}U_{x_{R-1}x_{R-2}}\cdots U_{x_1i}.
\]
Cada elección de índices intermedios determina un único camino de longitud
\(R\), y cada camino determina un único monomio ordenado. La conservación de
norma es consecuencia inmediata de \((U^R)^{\mathsf T}U^R=I\), o de su versión
unitaria.
\end{proof}

La fórmula \eqref{eq:suma-finita-caminos} es la identidad combinatoria que
subyace a la suma de historias. Su contenido HMT aparece al refinar cada
camino visible por las fibras de \(q_{\rm vis}\):
\begin{equation}
 \mathcal K_R(\widetilde f,\widetilde i)
 =\sum_{\substack{\gamma\in
        \mathsf P_{\TPK}(\widetilde i,\widetilde f)\\|\gamma|=R}}
 W(\gamma).
 \label{eq:kernel-genealogico-finito}
\end{equation}
En esa suma, dos rutas publicadas con iguales extremos siguen siendo términos
distintos mientras difieran en su genealogía.

\section{Realización unitaria sobre el registro genealógico \(27\times27\)}

La fuente histórica desarrolla una caminata cuántica sobre el toro
\[
 \Lambda_{27}=(\ZZ/27\ZZ)^2
\]
con cuatro orientaciones cardinales
\(\mathcal D=\{N,S,E,W\}\). Su espacio real de fase es
\[
 \mathcal H_{27}^{\RR}
 =\ell^2(\Lambda_{27})\otimes\RR^4\otimes E,
\]
donde \((E,J_E)\) realiza internamente la estructura compleja. El mezclador
de direcciones es la matriz de Walsh--Hadamard
\[
 C_4=\frac12
 \begin{pmatrix}
 1&1&1&1\\
 1&1&-1&-1\\
 1&-1&1&-1\\
 1&-1&-1&1
 \end{pmatrix},
 \qquad C_4^{\mathsf T}C_4=I_4.
\]
Una acción elemental orientada \(s_d\in\mathcal L_S\) define la puerta
real de fase
\[
 D_J=\operatorname{diag}
 \left(
  e^{-J_Es_N/\hbar_{\rm int}},
  e^{-J_Es_S/\hbar_{\rm int}},
  e^{-J_Es_E/\hbar_{\rm int}},
  e^{-J_Es_W/\hbar_{\rm int}}
 \right).
\]
El shift condicionado transporta cada componente en su dirección:
\begin{align*}
 (S\psi)(x,y,N)&=\psi(x,y-1,N),&
 (S\psi)(x,y,S)&=\psi(x,y+1,S),\\
 (S\psi)(x,y,E)&=\psi(x-1,y,E),&
 (S\psi)(x,y,W)&=\psi(x+1,y,W),
\end{align*}
con coordenadas módulo \(27\). El paso elemental y la composición
dodecafásica son
\begin{equation}
 U_{\rm micro}=S D_J C_4,
 \qquad
 U_{12}=U_{\rm micro}^{12}.
 \label{eq:unitario-ledger-27}
\end{equation}

\begin{proposition}[Unitariedad y suma exacta en el ledger]
El operador \(U_{\rm micro}\) es ortogonal sobre la realificación y unitario
en la carta compleja. Para todo \(n\), cada entrada de
\(U_{\rm micro}^n\) es la suma ordenada de los pesos de todos los caminos de
\(n\) pasos que unen sus estados extremos.
\end{proposition}

\begin{proof}
El shift es una permutación, \(C_4\) es ortogonal y cada bloque
\(e^{-J_Es_d/\hbar_{\rm int}}\) es una rotación. Su producto es ortogonal.
La segunda afirmación es el
\cref{thm:suma-finita-caminos} aplicado a la matriz por bloques
\(U_{\rm micro}\).
\end{proof}

La simetría traslacional reduce el problema espectral a bloques de tamaño
cuatro. Para
\[
 k_x=\frac{2\pi m}{27},\qquad
 k_y=\frac{2\pi n}{27},
 \qquad 0\le m,n<27,
\]
el shift se diagonaliza como
\[
 S(k_x,k_y)=\operatorname{diag}
 \bigl(e^{J_Ek_y},e^{-J_Ek_y},e^{J_Ek_x},e^{-J_Ek_x}\bigr),
\]
y
\begin{equation}
 U_{\rm micro}(k_x,k_y)=S(k_x,k_y)D_JC_4,
 \qquad
 U_{12}(k_x,k_y)=U_{\rm micro}(k_x,k_y)^{12}.
 \label{eq:bloque-fourier-ledger-27}
\end{equation}
El espectro global es la unión de los espectros de los \(27^2\) bloques.
La expresión combina dos representaciones de la misma dinámica: la potencia
matricial enumera caminos y la transformada de Fourier publica
cuasi-energías.

En cada bloque unitario puede elegirse un logaritmo espectral y definir
\[
 H_{\rm eff}(k_x,k_y)
 =\frac{\hbar_{\rm int}}{J_E\,\Delta t}
   \log U_{12}(k_x,k_y),
\]
entendido mediante cálculo funcional en la carta compleja o su realificación.
La elección de rama fija la zona de cuasi-energía. El Hamiltoniano efectivo
surge así del propagador ya construido; su exponencial recupera
\(U_{12}\) en cada modo.

\section{Carácter real de acción}

Sea \(\mathcal A_*:\mathsf P_{\TPK}\to\RR_{\geq0}\) el cociclo aditivo de
acción de la fuente HMT--MD. El calibre es
\[
 \lambda_*
 =1-\frac5{468}-\frac3{11}\Delta_4
 =0.9892859130191415\ldots\in(0,1).
\]
La cláusula local \(\mathsf H_{\rm loc}\) asigna coste unitario a un giro,
coste \(\lambda_*\) a un cambio de hoja y suma ambos cuando concurren en la
misma arista. Bajo esa cláusula,
\[
 \mathcal A_*(\gamma)
 =\#\operatorname{turns}(\gamma)
  +\lambda_*\#\operatorname{sheet\mbox{-}switches}(\gamma),
\qquad
 \mathcal A_*(\gamma_2\circ\gamma_1)
 =\mathcal A_*(\gamma_2)+\mathcal A_*(\gamma_1).
\]
La acción dimensional asociada es
\[
 S_{\rm int}(\gamma)=\hbar_{\rm int}\mathcal A_*(\gamma).
\]

En el módulo de fase real \((E,J_E)\) del capítulo anterior, con \(E\) un
espacio euclídeo real de dimensión finita, definimos
\begin{equation}
 \Phi_J(\gamma)
 :=\exp\!\left(-J_E\frac{S_{\rm int}(\gamma)}{\hbar_{\rm int}}\right)
 =\cos\mathcal A_*(\gamma)\,I
  -\sin\mathcal A_*(\gamma)\,J_E.
 \label{eq:caracter-real-accion}
\end{equation}

\begin{proposition}[Multiplicatividad de la fase interna]
\label{prop:multiplicatividad-fase-interna}
El mapa \(\Phi_J\) es un funtor de \(\mathsf P_{\TPK}\) al grupo
ortogonal generado por \(J_E\):
\[
 \Phi_J(\id_x)=I,
 \qquad
 \Phi_J(\gamma_2\circ\gamma_1)
 =\Phi_J(\gamma_2)\Phi_J(\gamma_1).
\]
\end{proposition}

\begin{proof}
La primera igualdad usa \(\mathcal A_*(\id_x)=0\). Para la segunda, la
aditividad de \(\mathcal A_*\) y la conmutación de múltiplos del mismo
endomorfismo \(J_E\) dan
\[
 e^{-J_E(\mathcal A_*(\gamma_2)+\mathcal A_*(\gamma_1))}
 =e^{-J_E\mathcal A_*(\gamma_2)}
  e^{-J_E\mathcal A_*(\gamma_1)}.
\]
\end{proof}

Sea ahora \(T(e)\in\End(E)\) el transporte acotado de fibra de una arista y
supóngase
que preserva \(J_E\). El peso
\[
 W_J(\gamma)
 =\prod_{e\in\gamma}^{\longleftarrow}
   T(e)\exp\!\left(-J_E\frac{s(e)}{\hbar_{\rm int}}\right)
\]
realiza enteramente sobre \(E\) la convención hamiltoniana
\(\exp(-iS_H[\gamma]/\hbar)\). La convención lagrangiana habitual de Feynman
\(\exp(+iS_L[\gamma]/\hbar)\) se obtiene tomando \(S_L=-S_H\), o,
equivalentemente, introduciendo una orientación
\(\varepsilon\in\{-1,+1\}\) en
\(\exp(\varepsilon J_ES/\hbar_{\rm int})\). El símbolo escalar \(i\) es la
carta compleja de la operación algebraica realizada por \(J_E\). El coste
positivo \(\mathcal A_*\) se convierte en acción física al aplicar el mapa de
realización correspondiente.

\section{Composición de propagadores sobre rutas concatenadas y ley de amplitudes}

La genealogía conserva la ley de composición que hace dinámico al kernel.

\begin{proposition}[Convolución y corte temporal]
\label{prop:chapman-kolmogorov-genealogico}
Supóngase que \(\widetilde X\) es finito; más generalmente, que la familia de
productos siguiente es absolutamente sumable. Si toda ruta de longitud
\(R+S\) posee un único corte tras \(R\) pasos, entonces
\begin{equation}
 \mathcal K_{R+S}(\widetilde f,\widetilde i)
 =\sum_{\widetilde x\in\widetilde X}
   \mathcal K_S(\widetilde f,\widetilde x)
   \mathcal K_R(\widetilde x,\widetilde i).
 \label{eq:convolucion-kernel}
\end{equation}
La identidad sigue siendo válida con pesos no conmutativos si se conserva el
orden causal de los factores.
\end{proposition}

\begin{proof}
Particionamos el conjunto de rutas por el estado enriquecido
\(\widetilde x\) alcanzado en el corte.
La factorización functorial del peso transforma la suma sobre rutas completas
en el producto de las dos sumas parciales, con el segmento posterior a la
izquierda.
\end{proof}

El corte sobre el estado visible sería válido sólo si peso y composición
descendieran fielmente por \(q_{\rm vis}\). En presencia de memoria, dos
tramos con el mismo visible pueden pertenecer a fibras incompatibles; por eso
la convolución se efectúa sobre \(\widetilde X\).

Esta ecuación distingue la suma de caminos de una colección de coincidencias
numéricas. El propagador posee semigrupo, extremos, composición y una ley de
refinamiento.

\section{Límite cilíndrico del sistema finito y construcción de la integral correspondiente}

Una integral funcional requiere controlar el espacio de rutas y su
refinamiento. Para extremos enriquecidos \(\widetilde i,\widetilde f\), sean
\(\mathcal P_n(\widetilde i,\widetilde f)\) los conjuntos finitos de caminos
admisibles a resolución \(n\), con restricciones sobreyectivas, y definamos
\[
 \Omega_{fi}^{\rm path}
 :=\varprojlim_n\mathcal P_n(\widetilde i,\widetilde f).
\]
Es el espacio compacto de genealogías de caminos con esos datos de frontera.
Sea \(\{\Pi_n\}_{n\geq0}\) una sucesión de particiones cilíndricas de
\(\Omega_{fi}^{\rm path}\), con \(\Pi_{n+1}\) refinando \(\Pi_n\). Para cada
cilindro \(C\in\Pi_n\), elijamos un representante
\(\gamma_C\), un peso proyectivo \(\mu_n(C)\) y una acción aproximada
\(S_n(C)\).

Definimos el operador simple
\begin{equation}
 \mathcal I_n(F)
 :=\sum_{C\in\Pi_n}
   \mu_n(C)
   F(\gamma_C)
   \exp\!\left(-J_E\frac{S_n(C)}{\hbar_{\rm int}}\right).
 \label{eq:integral-cilindrica-n}
\end{equation}

\begin{theorem}[Límite genealógico de integrales cilíndricas]
\label{thm:integral-cilindrica-limite}
Supóngase que:
\begin{enumerate}
 \item las masas cilíndricas son proyectivamente compatibles;
 \item \(F:\Omega_{fi}^{\rm path}\to\End(E)\) es continua y acotada;
 \item existe una función continua
       \(S:\Omega_{fi}^{\rm path}\to\mathcal L_S\) tal que
       \[
        \max_{C\in\Pi_n}\sup_{\gamma\in C}
        \left\|
          \frac{S_n(C)-S(\gamma)}{\hbar_{\rm int}}
        \right\|\longrightarrow0;
       \]
 \item la malla de las particiones,
       \(\operatorname{mesh}(\Pi_n)
        =\max_{C\in\Pi_n}\diam_{d_\vartheta}C\),
       tiende a cero.
\end{enumerate}
Como \(E\) es finito-dimensional, \(\End(E)\), provisto de la norma de
operador, es un espacio de Banach; la integral que sigue es por tanto una
integral de Bochner bien definida \cite{Bochner1955}.
Entonces \(\mathcal I_n(F)\) converge en norma de operador y su límite no
depende de los representantes. Escribimos
\begin{equation}
 \int_{\Omega_{fi}^{\rm path}}
 F(\gamma)
 \exp\!\left(-J_E\frac{S(\gamma)}{\hbar_{\rm int}}\right)
 \dd\mu(\gamma)
 :=\lim_{n\to\infty}\mathcal I_n(F).
 \label{eq:integral-genealogica}
\end{equation}
\end{theorem}

\begin{proof}
La compatibilidad proyectiva construye la medida de cilindros \(\mu\). La
función
\[
 G(\gamma)=F(\gamma)
 \exp\!\left(-J_E S(\gamma)/\hbar_{\rm int}\right)
\]
es continua sobre el compacto \(\Omega_{fi}^{\rm path}\), luego uniformemente
continua y acotada. Sea
\[
 G_n(C)=F(\gamma_C)
 \exp\!\left(-J_E S_n(C)/\hbar_{\rm int}\right).
\]
Como \(J_E\) es una estructura ortogonal,
\[
 \|e^{-J_Ea}-e^{-J_Eb}\|\leq |a-b|.
\]
La tercera hipótesis implica que
\(G_n(C)-G(\gamma_C)\) tiende uniformemente a cero. La cuarta hipótesis y la
uniforme continuidad de \(G\) hacen que las sumas con \(G(\gamma_C)\) sean
sumas de Riemann convergentes e independientes de los representantes. El
límite coincide con la integral de Bochner de \(G\) respecto de \(\mu\).
Más precisamente, si
\[
 \epsilon_n
 =\max_{C\in\Pi_n}\sup_{\gamma\in C}
   \left\|
    \frac{S_n(C)-S(\gamma)}{\hbar_{\rm int}}
   \right\|,
\]
entonces
\[
 \left\|\mathcal I_n(F)-\int G\,\dd\mu\right\|
 \leq
 \omega_F\!\bigl(\operatorname{mesh}(\Pi_n)\bigr)
 +\|F\|_\infty\epsilon_n,
\]
pues la integral se descompone sobre los cilindros y la masa total es uno.
Ambos términos tienden a cero por hipótesis.
\end{proof}

\begin{corollary}[Convergencia de una familia de kernels discretos]
\label{cor:convergencia-kernels-discretos}
Además de las hipótesis del
\cref{thm:integral-cilindrica-limite}, supóngase que existe una sucesión
cofinal de profundidades \(R_n\) y funciones simples \(F_{fi,n}\) tales que
\[
 \mathcal K_{R_n}(\widetilde f,\widetilde i)
 =\mathcal I_n(F_{fi,n}),
 \qquad
 F_{fi,n}\longrightarrow F_{fi}
\]
uniformemente. Entonces los kernels
\(\mathcal K_{R_n}(\widetilde f,\widetilde i)\) convergen a la integral
genealógica de \(F_{fi}\).
\end{corollary}

\begin{proof}
La diferencia entre \(\mathcal I_n(F_{fi,n})\) y
\(\mathcal I_n(F_{fi})\) está acotada por
\(\|F_{fi,n}-F_{fi}\|_\infty\), pues \(\mu_n\) tiene masa total uno y la
fase es ortogonal. El segundo término converge por el
\cref{thm:integral-cilindrica-limite}.
\end{proof}

El teorema produce una integral rigurosa sobre el espacio genealógico
compacto. El corolario explicita el puente que identifica cada kernel
matricial con una suma cilíndrica normalizada. La integral oscilatoria de
Feynman en tiempo real aparece al elegir
una representación física y un refinamiento cuya acción discreta aproxime la
acción del sistema. Esa identificación exige demostrar la convergencia del
propagador correspondiente; no queda concedida por la mera existencia de
\(\mu\).

\section{Integral visible e integral condicionada por memoria}

Un lector real de caminos
\(\Val_{\rm path}:\Omega_{fi}^{\rm path}\to K\subset\RR\) induce la medida
visible \(\nu=(\Val_{\rm path})_*\mu\). Para una función que factoriza por el
valor, \(F=\widehat F\circ\Val_{\rm path}\),
\[
 \int_{\Omega_{fi}^{\rm path}}F\,\dd\mu
 =\int_K\widehat F(x)\,\dd\nu(x).
\]
Sin embargo, la fase o el transporte pueden no factorizar por
\(\Val_{\rm path}\). La
desintegración de \(\mu\) sobre las fibras del valor expresa la diferencia:
\begin{equation}
 \int_{\Omega_{fi}^{\rm path}}G(\gamma)\,\dd\mu(\gamma)
 =\int_K
   \left[
    \int_{\Val_{\rm path}^{-1}(x)}G(\gamma)\,\dd\mu_x(\gamma)
   \right]\dd\nu(x).
 \label{eq:desintegracion-caminos}
\end{equation}
La integral interior conserva las historias que publican la misma coordenada.
Ésta es la forma matemática de integrar respecto de los valores sin reducir el
número a su cifra.

\section{Acción, interferencia y persistencia}

La suma genealógica mantiene simultáneamente tres operaciones:
\[
 \text{admisibilidad de ruta}
 \quad\longrightarrow\quad
 \text{peso/acción}
 \quad\longrightarrow\quad
 \text{interferencia en la representación}.
\]
La primera pertenece al TPK y a su ley de supervivencia. La segunda pertenece
al cociclo y al transporte dimensional. La tercera pertenece al módulo de
fase y al lector físico. Ninguna magnitud observada selecciona
retrospectivamente la ruta que la produce.

\begin{resultbox}
La suma finita de Feynman se obtiene como representación lineal exacta de la
categoría enriquecida de rutas. La estructura compleja puede realizarse por un
endomorfismo real \(J_E^2=-I\). Bajo compatibilidad de medidas, contracción de
cilindros y convergencia uniforme de acciones, las sumas cilíndricas definen
una integral genealógica. Los kernels finitos convergen a ella cuando,
además, satisfacen el puente del
\cref{cor:convergencia-kernels-discretos}.
\end{resultbox}

\begin{statusbox}
La expansión matricial finita, el cociclo de acción y las realizaciones reales
de la estructura compleja son resultados recuperados. Su ensamblaje en
\eqref{eq:caracter-real-accion} y \eqref{eq:kernel-genealogico-finito} es una
formalización nueva exacta. El teorema \ref{thm:integral-cilindrica-limite} es
exacto bajo sus premisas. La equivalencia con un modelo cuántico continuo
concreto requiere verificar esas premisas para el Hamiltoniano y la acción de
ese modelo.
\end{statusbox}

Las premisas del teorema identifican exactamente la estructura que pasa del
ledger finito al espacio funcional: compatibilidad proyectiva de los pesos,
contracción de la malla, convergencia uniforme de la acción y correspondencia
entre kernels y sumas cilíndricas. En la realización
\eqref{eq:unitario-ledger-27}, la composición y la unitariedad son identidades
algebraicas; al refinar el ledger, las restantes premisas determinan la
convergencia hacia el propagador continuo elegido.
